\section{Introduction}

La direction générale de MEDICA-SERVICES nous a mandatés pour évaluer la maîtrise de sa sécurité des systèmes d'information. L'entreprise héberge des applications de soins pour des établissements privés, et une défaillance de son SI peut toucher la prise en charge de patients chez ses clients. La question posée est double. Quel niveau de rigueur l'entreprise devrait-elle viser, et quel niveau ses processus démontrent-ils réellement ?

Nous appliquons la méthode de maturité SSI publiée par le SGDN et la DCSSI en novembre 2007, utilisée ici comme référence pédagogique historique et non comme certification ou preuve de conformité actuelle. Ses douze questions donnent le niveau adéquat et ses sept processus donnent les niveaux effectifs. L'audit porte uniquement sur les pièces \piece{P01} à \piece{P17} et \piece{T01} à \piece{T10}. Chaque conclusion cite sa pièce et distingue le fait, la déclaration, l'interprétation et l'information manquante. Aucun scan ni aucune connexion n'a été réalisé sur une infrastructure réelle.

Le rapport présente d'abord l'entreprise et le périmètre, puis le niveau adéquat et la maturité effective. Il expose ensuite l'audit technique et sa corrélation avec la gouvernance, avant de proposer les risques, l'arbitrage budgétaire et la feuille de route. L'organisation de notre travail et la conclusion le referment. Les tableaux détaillés figurent en annexe.
